\subsection{Tensor realization and the exact arithmetic comparison gate}
The fixed-parameter action extends through all tensor degrees.
The stronger requirement that finite Weil centralizers recover the
entire nonstandard endomorphism system has different consequences.
Rational dependence on the quantum parameter makes the extension
comparison especially restrictive.

\begin{finding}[\fid{F48}: faithful tensor coefficients and full-centralizer obstructions]\label{f:tensorweilgate}
In the actual two-row Kronecker setting let $X=V\otimes W$,
with $\dim V=\dim W=2$, and fix a coefficient parameter at which
the indicated source maps are defined.
The finite-dimensional nonstandard polynomial-comodule category
has an exact faithful strong monoidal realization by finite
Weil coefficient objects, all pure of weight zero after an
arithmetic coefficient specialization. In IV's actual Hopf
setting this realization preserves the existing rigid duals.
It retains every named tensor action and every specified
source map. Nontrivial finite-order tensor Frobenii also exist.

The realization cannot be fully faithful into the category of
finite-dimensional Weil representations on a point.
The actual polynomial endomorphism algebras already satisfy
\[
 \End(X)=E,\qquad
 \End(X^{\otimes2})=E\,1\oplus E\,P,
\]
where $P$ is the rank-ten nonstandard symmetric projector.

There are additional arithmetic restrictions independent of
that categorical fullness condition. Let $b>1$ and let
$F(v)\in GL_N(\C(v))$ be regular and invertible at $b$ and
at all sufficiently large $b^n$. If
\begin{equation}\label{eq:tensorweilspectral}
 \operatorname{Spec}F(b^n)
   =\operatorname{Spec}(F(b)^n)
\end{equation}
as eigenvalue multisets for every sufficiently large integer $n$,
then, for finitely many nonnegative integers $r_m$,
\[
 \operatorname{Tr}F(v)=\sum_m r_m v^m,\qquad
 \det(TI-F(v))=\prod_m(T-v^m)^{r_m}.
\]
The $r_m$ are the actual positive multiplicities of the
eigenvalues that occur.

Consequently no such rational family of one fixed integral
pure weight can recover the actual nonsimple semisimple
rank-three algebra at $b>1$ as its full Frobenius centralizer,
at any auxiliary rank. If $F(v)$ is also regular at $v=1$,
it cannot recover the full actual collision algebra there,
even with arbitrary mixed weights.

Under the stronger equality of matrices
$F(b^n)=F(b)^n$ in one fixed rational trivialization,
there is a constant invertible matrix $S$ and integers $m_j$
such that
\[
 F(v)=S\operatorname{diag}(v^{m_1},\ldots,v^{m_N})S^{-1}.
\]
These restrictions do not exclude the faithful actions as
specified subalgebras supplied by the first part.
\end{finding}

\begin{proof}
\par\smallskip\noindent\textbf{Actual source data and conventions.}\quad
Write $B$ for IV's nonstandard matrix coordinate bialgebra.
Its defining relation is
\[
 P(u\otimes u)=(u\otimes u)P;
\]
see IV, Definition~8.1, equation~(61), \S8.1, p.~39.
In this Kronecker setting, Proposition~10.3, \S10.1, p.~56,
constructs the Hopf algebra after adjoining the inverse
nonstandard determinant. Proposition~10.5, \S10.2, p.~57,
gives its Hopf map to the product of the ordinary quantum
$GL(V)$ and $GL(W)$ coordinate algebras; the determinant
maps to $D_v(V)^2D_v(W)^2$.
Theorem~12.1, p.~68, identifies the degree-$r$ coordinate
dual with the commutant of the adjacent projectors
\cite{IV}.

The remark immediately after that theorem corrects its
displayed $X$ to $X^*$ relative to the conventions in
\S\S8--10, and its proof uses the dual consistently.
Here one matrix corepresentation and its induced tensor
action are fixed throughout. Passing to its dual or
transposed presentation preserves the endomorphism
dimensions and all obstruction arguments below; the
source projector is symmetric in its chosen coordinates.

The Hopf assertion just cited is specific to IV.
VII, \S2, explicitly allows the general plethysm determinant
to vanish. Its bialgebra category still has tensor products;
no general VII antipode or rigid duals are inferred here.
At an algebraic specialization we use only the coefficient
maps actually defined there, without importing a generic
compact-real-form theorem into an unverified specialization.

Put
\[
 \delta=v+v^{-1},\qquad D=\delta^2,\qquad
 P=P_+^V\otimes P_+^W+P_-^V\otimes P_-^W.
\]
The ordinary symmetric and alternating projectors have ranks
$3$ and $1$. Thus $P$ has rank $3\cdot3+1\cdot1=10$,
and $1-P$ has rank $6$. On $X^{\otimes r}$ the actual
right generators are $DP_i$ on consecutive tensor factors.
These assertions hold in particular at positive real
parameters and at $v=1$.

\par\smallskip\noindent\textbf{The faithful tensor realization.}\quad
Let $\omega$ be the ordinary forgetful functor on finite
$B$-comodules. Send
\[
 N\longmapsto(\omega N,1)
\]
to finite-dimensional vector spaces with an invertible
linear operator. Kernels and cokernels of comodule maps
are computed on underlying vector spaces, so this functor
is exact and faithful. The bialgebra multiplication defines
the source tensor product; its underlying tensor product
is the ordinary one. Since $1\otimes1=1$, the ordinary
associativity and unit comparisons are Weil isomorphisms
with the required coherence.
In a rigid Hopf category the ordinary underlying dual,
evaluation and coevaluation are unchanged, and the dual
of the identity is the identity. Thus this also retains
the specified IV dual maps and determinant inversion.
After an arithmetic coefficient embedding all eigenvalues
are one, giving purity of weight zero and a continuous
point-Galois representation.

The right generators on each image remain exactly the
original matrices $DP_i$. Tensoring two consecutive blocks
gives their actual block inclusion, and all original
relations remain matrix equalities.
Every source comodule map remains the same linear map.
In addition, ordinary left operators that are not comodule
intertwiners, including the named operators of \fid{F43},
are still Weil endomorphisms of these coefficient objects,
because they commute with identity Frobenius.
This is retention of the actual operations, not a claim
that arbitrary linear maps become source comodule maps.

One may assign a fixed pure unit scalar $c$ to $X$ and
$c^r$ to homogeneous degree $r$, with inverse on dual
degree, when that grading of all maps is part of the
data. The identity construction requires no additional
grading qualification and is the unconditional version.

\par\smallskip\noindent\textbf{Nontrivial pure tensor characters.}\quad
Choose roots of unity $\zeta,\eta$ in the coefficient field.
The diagonal evaluations
\[
 A=\operatorname{diag}(\zeta,\zeta^{-1}),\qquad
 B'=\operatorname{diag}(\eta,\eta^{-1})
\]
are characters of the ordinary quantum $GL_2$ coordinate
algebras: off-diagonal generators go to zero, the remaining
relations hold for commuting diagonal scalars, and the
quantum determinants go to one. Compose with IV's actual
Hopf map to obtain a character $\chi$ of the nonstandard
Hopf algebra. On $X$ its matrix is $A\otimes B'$.
Directly in the FRT bialgebra the same character exists
because $(A\otimes B')^{\otimes2}$ commutes with $P$.

For a right comodule with coaction $\rho_N$, define
\[
 F_N=(1\otimes\chi)\rho_N.
\]
It is natural for all comodule maps. Multiplicativity of
$\chi$ gives
\[
 F_{N\otimes N'}=F_N\otimes F_{N'}.
\]
In the Hopf setting the inverse uses $\chi$ composed with
the antipode; on a dual the operator is the inverse
transpose. Choose a common order $L$ of $\zeta,\eta$.
Then $F_X^L=1$, and the same holds on every tensor product,
dual, stable subspace and quotient in the tensor category
they generate. In characteristic zero these operators
are semisimple with root-of-unity eigenvalues.
They are pure of weight zero under every embedding and
give continuous representations at every coefficient prime.

All actual right generators and source comodule maps
commute with these Frobenii. The ordinary left raising
and lowering operators have the precise covariance
\[
\begin{aligned}
 F_XE_VF_X^{-1}&=\zeta^2E_V,&
 F_XF_VF_X^{-1}&=\zeta^{-2}F_V,\\
 F_XE_WF_X^{-1}&=\eta^2E_W,&
 F_XF_WF_X^{-1}&=\eta^{-2}F_W.
\end{aligned}
\]
This follows first for the elementary upper and lower
matrices on each two-dimensional factor; the source
coproduct preserves the same equations on tensor powers.
Thus, for example, $E_V$ is a Weil map
$L_{\zeta^2}\otimes X\to X$, where $L_a$ is the rank-one
coefficient object with Frobenius $a$.
With identity Frobenius all these named operators are
ordinary invariant endomorphisms instead.
The character construction fixes the parameter; it does
not impose the extension-retuning rule.

\par\smallskip\noindent\textbf{Actual endomorphisms and failure of fullness.}\quad
In coordinate degree one there are no FRT relations.
The degree-one dual therefore acts as $\End_E(X)$,
whose commutant is $E$. In degree two the sole relations
are the entries of $PU-UP$, where $U$ is the matrix
of the universal tensor-square corepresentation.
Their annihilator identifies the acting coordinate dual
with
\[
 \{a\in\End_E(X^{\otimes2}):aP=Pa\}
   =\End(\operatorname{im}P)\oplus\End(\ker P).
\]
Taking the commutant of these two full matrix blocks gives
\[
 \End_B(X)=E,\qquad
 \End_B(X^{\otimes2})=E\,1\oplus E\,P.
\]
The blocks are nonzero of ranks ten and six.
This is also the degree-two instance of the source
proof of Theorem~12.1.

For the IV Hopf category, the facts needed for the same
obstruction are $\End_H(X)=E$ and the independent
endomorphisms $1,P$ of $X^{\otimes2}$.
The first follows by ordinary quantum-product restriction:
the elementary raising and lowering matrices on the two
fundamental factors generate all four-by-four matrix
units on $X$, so a commuting endomorphism is scalar.
The second is the defining FRT relation after determinant
inversion. No assertion that arbitrary localization
preserves every polynomial Hom space is needed.

Suppose a fully faithful $E$-linear strong monoidal
functor sent $X$ to a nonzero Weil object $(Y,F)$.
Fullness in degree one gives $\End_F(Y)=E$.
Since $F$ itself belongs to this algebra, $F=cI$.
Its full centralizer is then $\End_E(Y)$, so $\dim Y=1$.
Strong monoidality makes the image of $X^{\otimes2}$
one-dimensional as well, contradicting its independent
endomorphisms $1,P$. This proves the fullness obstruction.

Equivalently, auxiliary objects $H_r$ with
$H_{r+s}\simeq H_r\otimes H_s$ cannot identify the
degree-one and degree-two nonstandard algebras with
their entire Frobenius centralizers.
The rank-six $H$ of \fid{F46} is an auxiliary degree-three
presentation object and has not been asserted to be
such a tensor power.
This obstruction already occurs in the generic
semisimple setting; it is distinct from the collision
radical obstruction in \fid{F47}.

\par\smallskip\noindent\textbf{A scalar extension-law classification.}\quad
For later comparison, let $f\in\C(v)^\times$ be regular
and nonzero at $b>1$ and suppose
$f(b^n)=f(b)^n$ for all sufficiently large $n$.
Write its leading Laurent term at infinity as
\[
 f(v)=a v^m(1+O(v^{-1})),\qquad a\ne0,\quad m\in\Z.
\]
The ratio of consecutive evaluations tends to $b^m$,
but is identically $f(b)$ by the extension law.
Thus $f(b)=b^m$. The rational function $f(v)-v^m$
vanishes at infinitely many $b^n$ and is zero.
Conversely every $v^m$ satisfies the law.
In particular a nontrivial constant phase is not allowed
by this all-extension equation.
For $b=Q^{1/2}$ its weight is $m$, so weight zero forces
$f=1$. The inverse parameter convention is obtained by
replacing $v$ by $v^{-1}$.
The relation $f(v^{-1})=f(v)^{-1}$ already holds for
these monomials. This statement concerns a rational
labeled branch, not an unspecified multivalued
algebraic continuation.

\par\smallskip\noindent\textbf{Unlabeled rational spectra.}\quad
We now prove the stronger matrix-spectral statement
without labeling branches or assuming diagonalizability.
Put $f(v)=\operatorname{Tr}F(v)$.
Under \eqref{eq:tensorweilspectral},
\[
 f(b^n)=\sum_\alpha m_\alpha\alpha^n,
\]
where the distinct nonzero $\alpha$ are eigenvalues of
$F(b)$ and $m_\alpha>0$ are their actual multiplicities.
This finite power sum satisfies a constant-coefficient
recurrence
\[
 \sum_{j=0}^s c_j f(b^{n+j})=0,\qquad c_0\ne0.
\]
The constant term is nonzero because all roots are nonzero.
The rational function
\[
 G(v)=\sum_{j=0}^s c_j f(b^jv)
\]
vanishes at infinitely many $b^n$, so it is identically
zero. If $f$ had a nonzero finite pole $\rho$ of maximal
modulus, a pole of $f(b^jv)$ for $j\ge1$ would have modulus
at most $|\rho|/b^j<|\rho|$.
Only $c_0f(v)$ could have a pole at $\rho$, a contradiction.
Thus $f$ has poles only at zero and infinity and is a
Laurent polynomial. This is the rational dilation
argument of \fid{F26}, reproved here in the needed form.

Write $f(v)=\sum_m a_m v^m$. Evaluating gives
\[
 \sum_m a_m(b^m)^n=\sum_\alpha m_\alpha\alpha^n.
\]
Distinct nonzero finite power roots are linearly
independent: any zero relation at sufficiently many
consecutive integers gives an invertible Vandermonde
matrix, times a nonzero diagonal matrix of root powers.
Therefore every $\alpha$ is $b^m$, and $a_m$ is its
positive integer multiplicity $r_m$.
This proves the Laurent trace assertion.

At $v=b^n$ the characteristic polynomial is now
$\prod_m(T-b^{mn})^{r_m}$.
Each rational coefficient agrees at infinitely many
points with the corresponding coefficient of
$\prod_m(T-v^m)^{r_m}$, proving the family identity.
Virtual alternating traces with cancelled roots would
not justify this inference about invisible eigenvalues;
the hypothesis here is the full matrix spectrum.

\par\smallskip\noindent\textbf{One pure weight and the generic actual algebra.}\quad
Suppose $b=Q^{1/2}$ and $F(b)$ is pure of one fixed
integral weight $w$. Normalize
\[
 F_0(v)=v^{-w}F(v).
\]
It is rational, has the same full centralizer, and still
satisfies \eqref{eq:tensorweilspectral}, since the
normalizing factor at $b^n$ is $(b^{-w})^n$.
Every eigenvalue of $F_0(b)$ has modulus one.
The spectral classification forces them all to be one,
so $N_0=F_0(b)-I$ is nilpotent.

At $b>1$, the actual rank-three algebra of \fid{F19} is
\[
 A_b\simeq E\times E\times M_2(E)\times M_2(E)
\]
over its splitting coefficient field.
Indeed $D=(b+b^{-1})^2>4$ and
$\Delta=(D-1)(D-4)/D^2\ne0$, so its center branches split.
This algebra is semisimple and not simple.
If it were the full centralizer of $F_0(b)$, the
nilpotent $N_0$ would be central in that semisimple algebra,
because every centralizer element commutes with it.
A semisimple finite algebra has reduced center, so
$N_0=0$. But the centralizer of $I$ is $M_N(E)$,
which is simple, a contradiction.
The argument permits every auxiliary rank and every
Jordan form. For $v=Q^{-1/2}$ the normalization is
$v^wF(v)$, or equivalently the preceding proof after
inverting the parameter.

\par\smallskip\noindent\textbf{Mixed weights and the retained collision.}\quad
We next prove the center fact needed at $v=1$.
For a nilpotent operator $N_0$ of nilpotency index $s$
on a nonzero finite-dimensional $k$-space,
\begin{equation}\label{eq:nilpotentcentralizercenter}
 Z(\End_{N_0}(Y))=k[N_0]\simeq k[t]/(t^s).
\end{equation}
Regard $Y$ as the $k[t]$-module
$\bigoplus_i k[t]/(t^{r_i})$, with $\max r_i=s$.
A central endomorphism commutes with the projections
onto the cyclic summands, so is block diagonal.
Its restriction to each summand is multiplication by
a polynomial modulo $t^{r_i}$.
Choose a largest summand $k[t]/(t^s)$ and let $f$
be that restriction. The quotient map from it to each
smaller summand, sending $1$ to $1$, is an intertwiner.
Commuting with these maps forces every restriction to
be $f$ modulo the appropriate power of $t$.
This also handles multiple largest summands.
Conversely every polynomial in $N_0$ is central.
The largest summand makes the kernel exactly $(t^s)$,
proving \eqref{eq:nilpotentcentralizercenter}.
Its center is local, over any field.

If the rational family $F(v)$ is regular at $v=1$,
the proved characteristic-polynomial identity gives
\[
 \det(TI-F(1))=(T-1)^N.
\]
Thus $F(1)-I$ is nilpotent and its full centralizer has
local center. The actual collision algebra, however, is
\[
 A_1=k\times k\times M_2(k[\epsilon]/(\epsilon^2)),
\qquad
 Z(A_1)=k\times k\times k[\epsilon]/(\epsilon^2).
\]
Its center has three nonzero primitive central idempotents
and is not local. Hence it cannot be the full centralizer
of $F(1)$.
No purity assumption was used in this collision argument.
A rational change of basis giving a regular matrix model
does not change its characteristic-polynomial identity
or its local-center conclusion.
The assumptions are a finite fixed-rank family, regularity
at the collision, and recovery of the full special
centralizer; a nonflat coherent object of different
special rank or a specified proper subalgebra is outside
this assertion.

\par\smallskip\noindent\textbf{The stronger equality of matrices.}\quad
Finally suppose $F(b^n)=F(b)^n$ in a fixed rational
trivialization. Choose a constant complex Jordan basis
$S$ for $F(b)$ and set $G(v)=S^{-1}F(v)S$.
Every diagonal entry satisfies
$g_j(b^n)=\alpha_j^n$, so the scalar argument above gives
$g_j(v)=v^{m_j}$ and $\alpha_j=b^{m_j}$.

If a Jordan block had length at least two, its first
superdiagonal entry $h(v)$ would satisfy
\[
 h(b^n)=n\alpha^{n-1}.
\]
The ratio of consecutive values tends to
$\alpha=b^m$, forcing the leading Laurent exponent
of $h$ to be $m$.
But $h(b^n)/b^{mn}$ tends to its finite leading
coefficient, whereas the displayed value is
$n/\alpha$, unbounded. This is impossible.
All Jordan blocks have length one.
Every off-diagonal entry now vanishes at infinitely
many $b^n$ and is zero rationally.
Consequently
\[
 F(v)=S\operatorname{diag}(v^{m_1},\ldots,v^{m_N})S^{-1}.
\]
Its converse is immediate. Pure weight zero forces
all $m_j=0$ and hence $F(v)=I$.
This final classification requires equality in a
fixed rational trivialization; spectral equality,
parameter-dependent conjugacy or independently chosen
isomorphisms do not imply that stronger hypothesis.
\end{proof}

\paragraph{What has been closed, and what still requires a comparison.}
There are exact finite pure coefficient tensor families
carrying the actual fixed-parameter actions and the IV
rigid dual maps. Thus absence of any pure tensor realization
is not the remaining issue.
Full recovery from finite point-Weil centralizers is a
stronger requirement and fails as above.
Rational all-extension spectral retuning with one fixed
pure weight already conflicts with the actual generic
full centralizer; allowing mixed weights still cannot
retain the full collision centralizer in a regular
fixed-rank family.
These statements allow subalgebras of larger centralizers,
additional comodule or geometric structure, nonrational
comparison data, and separately chosen word objects.
They do not supply the \fid{F24} opposite-sign output,
the required grading comparison, or VIII's prescribed
global integral and infinity lattices and basis selection.
The results concern finite matrix spectra with genuine
multiplicities, and do not bound eigenvalues invisible
in a virtual trace.

% Separate paragraph for Part III.
\paragraph{\fid{F48} in Grothendieck's language.}
The tensor realization starts from the actual bialgebra
coaction and its forgetful fiber functor. The precise
coaction and equivariant-module framework is SGA~3,
I, 4.7.2--3, with cote~90, scan~17 as the audited archive
anchor. The representation-ring and $\lambda$-operation
vocabulary for ordinary coefficient objects is in
SGA~6, V, \S\S2--3, with cote~133, scan~18 as the audited anchor;
it does not identify the nonstandard exterior relations
with the symmetric constraint of the Weil target.
Purity here uses the eigenvalue and complex conventions
of BBD~5.1.2 and~5.1.8.
The displayed low-degree commutants, rational recurrence
and pole argument, and cyclic-module center calculation
prove the specific additional restrictions.
These sources supply their frameworks, not an unpublished
nonstandard tensor or arithmetic comparison.
